\begin{frame}
    \begin{columns}[t]
        \begin{column}{0.48\textwidth}
            \begin{block}{Adomian decomposition method}
                The Adomian decomposition method, which accurately computes the series solution, is of great interest in applied sciences.
                The method provides the solution in a rapidly convergent series whose components are elegantly computed.

                A more difficult question to address is how this method can be modified to handle the concept of singular points.
                It is well known that the Adomian decomposition method usually begins by defining the equation in an operator form,
                considering the highest-order derivative in the problem.
                However, a slight modification is necessary to overcome the singular behavior.

            \end{block}
        \end{column}
        \begin{column}{0.48\textwidth}
            \begin{block}{Differential transformation method}
                Aquí nos centraremos en el método propuesto por el colombiano Cárdenas.
                En este aborda para $f(x)=e^{\pm x}$ y $f(x)=(x^2-C)^{\frac{3}{2}}$.
                Este método se basa en la misma lógica que la Transformadas de Laplace.
                Definimos el operador:
                $$Y(k)=\frac{1}{k!}[\frac{d^ky(x)}{dx^k}]_{x=x_0}.$$
                $y(x)=\sum_{k=0}^{\infty}Y(k)x^k$,junto con otros teoremas.
            \end{block}
        \end{column}
    \end{columns}
\end{frame}

% Frame 1: Planteamiento de la Serie de Potencias
\begin{frame}{Descomposición Modificada: Planteamiento}
    La solución puede obtenerse mediante el método de ``descomposición modificada''.
    Escribimos $T$ como una serie de potencias:

    \begin{equation}
        T=
        \sum_{n=0}^{\infty}
        a_{n} r^{n},\quad T^{m}=
        \sum_{n=0}^{\infty} A_{n} r^{n},
    \end{equation}

    donde $A_{n}=A_{n}(a_{0}, \dots, a_{n})$.

    Las condiciones iniciales imponen restricciones sobre los coeficientes:
    \begin{itemize}
        \item La condición $T(0)=1$ implica $\sum_{n=0}^{\infty} a_{n} r^{n}|_{r=0} = 1 \implies a_{0} = 1$.
        
        \item La condición $dT/dr|_{r=0} = 0$ implica $\sum_{n=0}^{\infty} (n+1)a_{n+1} r^{n}|_{r=0} = 0 \implies a_{1} = 0$.
    \end{itemize}
\end{frame}

% Frame 2: Operador Diferencial y Derivadas
\begin{frame}{Aplicación del Operador Diferencial}
    Calculamos las derivadas término a término para sustituir
    en el operador $L = r^{-2}(d/dr)r^{2}(d/dr)$:

    \begin{align*}
        \frac{dT}{dr}                               & = \sum_{n=1}^{\infty} n a_{n} r^{n-1} = \sum_{n=1}^{\infty} (n+1) a_{n+1} r^{n} \\
        r^{2}\frac{dT}{dr}                          & = \sum_{n=1}^{\infty} (n+1) a_{n+1} r^{n+2}                                     \\
        \frac{d}{dr}\left(r^{2}\frac{dT}{dr}\right) & = \sum_{n=0}^{\infty} (n+2)(n+3) a_{n+2} r^{n+1}
    \end{align*}

    Por lo tanto:
    \begin{equation}
        L T = r^{-2}\frac{d}{dr}r^{2}\frac{dT}{dr} = \sum_{n=0}^{\infty} (n+2)(n+3) a_{n+2} r^{n}.
    \end{equation}
\end{frame}

% Frame 3: Relación de Recurrencia
\begin{frame}{Relación de Recurrencia}
    Sustituyendo en la ecuación original $LT + \lambda T^m = 0$:

    \begin{equation}
        \sum_{n=0}^{\infty} (n+2)(n+3) a_{n+2} r^{n} + \lambda \sum_{n=0}^{\infty} A_{n} r^{n} = 0.
    \end{equation}

    Igualando coeficientes, obtenemos la relación:
    \begin{equation}
        (n+2)(n+3) a_{n+2} + \lambda A_{n} = 0.
    \end{equation}

    Esto nos permite determinar los coeficientes de la serie de Maclaurin para $T$:
    \begin{equation}
        a_{n+2} = -\frac{\lambda A_{n}}{(n+2)(n+3)}, \quad n \ge 0.
    \end{equation}

    Recordando que $a_{0}=1$ y $a_{1}=0$.
\end{frame}

% Frame 4: Polinomios de Adomian
\begin{frame}{Cálculo de los Polinomios de Adomian $A_n$}
    Los polinomios $A_n$ para no linealidades de potencia decimal se generan como:

    \begin{align*}
        A_{0} & = a_{0}^{m}                                                                                           \\
        A_{1} & = m a_{0}^{m-1} a_{1}                                                                                 \\
        A_{2} & = m a_{0}^{m-1} a_{2} + \frac{1}{2}m(m-1) a_{0}^{m-2} a_{1}^{2}                                       \\
        A_{3} & = m a_{0}^{m-1} a_{3} + m(m-1) a_{0}^{m-3} a_{1} a_{2} + \dots                                        \\
        A_{4} & = m a_{0}^{m-1} a_{4} + m(m-1) a_{0}^{m-2} \left[ \frac{1}{2} a_{2}^{2} + a_{1} a_{3} \right] + \dots
    \end{align*}

    Dado que $a_1 = 0$, todos los componentes de índice impar ($a_3, a_5, \dots$) y sus respectivos $A_n$ se anulan.
\end{frame}

% Frame 5: Solución Final
\begin{frame}{Solución de la Ecuación Lane-Emden}
    Simplificando con los términos nulos, la solución se expresa como:

    \begin{equation}
        T = 1 - \lambda \sum_{k=1}^{\infty} \frac{A_{2k-2}}{(2k)(2k+1)} r^{2k},
    \end{equation}

    donde los primeros coeficientes $A_n$ no nulos son:
    \begin{align*}
        A_{0} & = 1                                \\
        A_{2} & = -\frac{\lambda}{2 \cdot 3}       \\
        A_{4} & = -\frac{\lambda A_{2}}{4 \cdot 5}
    \end{align*}

    Este método de descomposición modificada permite obtener la solución sin necesidad de cambios de variable.

    NOTA:Hay un analogo para el MTD que es la tesis del alumno de pablo alzate.
    En este verifica las soluciones tanto en los índices m=0,1,2,3,4,5;Solo para esos números enteros.
    ¿Qué hago? Muestro el MTD así como mostré el Metodo de Adomiam?
\end{frame}